% =========================================================================
% Seguimiento de riesgos — 25-10-2026 (S9, tercer sprint de F2).
% Revisión por cadencia de los críticos (R1–R7 y R10) y de los significativos
% (R8, R9, R11, R14 y R15, que tocaban en S9); R12 por disparador (indicador
% saltado el 19-10). R13 (moderado) toca en S11. Sin plan nuevo: el único
% cambio es el estado de R12, que desde hoy se registra en la tabla del
% cap. 7 y aquí sin reemitir el plan (decisión del autor de este seguimiento).
% Auditado contra el dev-journal (21-10 y 25-10) y git. Documento
% independiente de la memoria; creado copiando TEMPLATE.tex (ver README.md
% de esta carpeta).
% Compilar DESDE ESTA CARPETA (2 pases, sin biber):
%   pdflatex --enable-installer -interaction=nonstopmode -file-line-error
%     -output-directory=build 2026-10-25-S9.tex
% =========================================================================
\documentclass[11pt,a4paper]{article}

% =========================================================================
% Preámbulo COMPARTIDO de los documentos independientes de riesgos
% (docs/thesis/plan-riesgos/ y docs/thesis/seguimiento-riesgos/).
%
% Estos documentos NO forman parte de main.tex (la memoria): son entregables
% sueltos en clase article, con el diseño heredado de los .tex de Gestión de
% Proyectos del autor, pero adaptados a la configuración global de la
% memoria (Carlito, márgenes 2,5 cm, parskip 6 pt, títulos de la guía).
% Un solo preámbulo para toda la familia: un ajuste de estilo aquí se
% propaga a todos los seguimientos futuros (uno por sprint hasta S37).
%
% RUTAS: los \input son relativos a la carpeta del documento, no a esta;
% compilar SIEMPRE desde la carpeta del .tex (plan-riesgos/ o
% seguimiento-riesgos/):
%   pdflatex --enable-installer -interaction=nonstopmode -file-line-error
%     -output-directory=build <documento>.tex        (2 pases; sin biber)
% =========================================================================

% --- Codificación (idéntico a packages.tex de la memoria) ---
\usepackage[utf8]{inputenc}
\usepackage[T1]{fontenc}

% --- Idioma: solo español (aquí no hay resumen trilingüe) ---
% es-noshorthands por coherencia con la memoria: acentos en UTF-8, sin
% caracteres activos (el conflicto con biblatex-apa no aplica aquí, pero
% mantener una única configuración de babel evita sorpresas al copiar texto
% entre la memoria y estos documentos).
\usepackage[spanish, es-tabla, es-noshorthands]{babel}

% --- Márgenes: 2,5 cm laterales e inferior, como la memoria ---
% El margen superior se reparte para alojar la cabecera enmarcada del
% diseño de referencia (2 cm de caja + separación) sin comerse el cuerpo:
% el texto arranca a 4 cm del borde y el ancho de texto sigue siendo 16 cm.
\usepackage[a4paper, left=2.5cm, right=2.5cm, bottom=2.5cm,
            top=1.2cm, includehead, headheight=2.2cm, headsep=0.6cm]{geometry}

% --- Tipografía: Carlito, la misma decisión que en style.tex ---
% (clon métrico libre de Calibri; pdflatex no puede usar la real).
% No se replica el arreglo de versalitas de style.tex: aquí no hay
% bibliografía APA y ningún texto usa \textsc.
\usepackage[sfdefault]{carlito}

% --- Espaciado: interlineado sencillo + 6 pt entre párrafos (guía) ---
\usepackage[skip=6pt]{parskip}

% --- Tablas ---
% longtable: el registro y el plan de respuesta superan una página.
% AVISO ESPERADO: "ignored error: Infinite glue shrinkage" en cada salto de
% página de una longtable (bug conocido e inofensivo, latex3/latex2e#1907;
% ver el comentario largo en packages.tex de la memoria).
\usepackage{array}
\usepackage{longtable}
\usepackage{tabularx}
% \RaggedRight (ragged2e) en vez de \raggedright en columnas muy estrechas:
% el \raggedright de LaTeX desactiva de facto la partición silábica y una
% palabra larga («Identificado») desborda una columna de 1,2 cm; la
% variante de ragged2e mantiene el guionado.
\usepackage{ragged2e}
% Las columnas X se componen como m{...}: el contenido de cada celda queda
% centrado verticalmente (aspecto de formulario del diseño de referencia).
\renewcommand{\tabularxcolumn}[1]{m{#1}}
% Tipos de columna del diseño de referencia:
%   L/C = ancho fijo alineado izquierda/centro; Y/Z = X alineada izq./centro.
% En columnas estrechas se prefiere Y (raggedright) a X justificada: la
% justificación en 3-4 cm abre ríos de blancos.
\newcolumntype{L}[1]{>{\raggedright\arraybackslash}m{#1}}
\newcolumntype{C}[1]{>{\centering\arraybackslash}m{#1}}
\newcolumntype{Y}{>{\raggedright\arraybackslash}X}
\newcolumntype{Z}{>{\centering\arraybackslash}X}
% Filas con aire en todas las tablas (el diseño de referencia lo fijaba
% tabla a tabla; aquí una sola vez).
\renewcommand{\arraystretch}{1.25}

% --- Utilidades ---
\usepackage{xcolor}    % \markfield (resaltado) y \todonote (macros.tex)
\usepackage{amssymb}   % \square, \blacksquare, \circlearrowright
\usepackage{xspace}    % lo necesita \bukit (macros.tex)
\usepackage{enumitem}
\setlist[itemize]{leftmargin=*, topsep=2pt, itemsep=2pt}
\setlist[enumerate]{leftmargin=*, topsep=2pt, itemsep=2pt}

% --- Títulos según la guía de estilo, como en style.tex de la memoria ---
% (sin \chapter: la clase article no lo tiene; el resto, idéntico).
\usepackage{titlesec}
\titleformat{\section}
  {\normalfont\bfseries\fontsize{12pt}{14.4pt}\selectfont\raggedright}
  {\thesection.}{0.5em}{\MakeUppercase}
\titleformat{\subsection}
  {\normalfont\bfseries\fontsize{12pt}{14.4pt}\selectfont}
  {\thesubsection.}{0.5em}{}
\titleformat{\subsubsection}
  {\normalfont\itshape\fontsize{12pt}{14.4pt}\selectfont}
  {\thesubsubsection.}{0.5em}{}

% --- Paginación dinámica (regla estricta del autor, 26-07-2026) ---
% PROHIBIDO el salto de página manual también en estos documentos: el
% espacio lo garantizan estos ganchos (o un \Needspace* local ante un bloque
% indivisible), nunca un \newpage. Los .tex de referencia sí usaban
% \newpage; es la única desviación deliberada respecto a ellos.
\usepackage{needspace}
\AddToHook{cmd/section/before}{\Needspace*{5\baselineskip}}
\AddToHook{cmd/subsection/before}{\Needspace*{5\baselineskip}}
\AddToHook{cmd/subsubsection/before}{\Needspace*{4\baselineskip}}
\widowpenalty=10000
\clubpenalty=10000
\displaywidowpenalty=10000
\raggedbottom
% Estos formularios componen mucho texto en columnas de 2-4 cm: se concede
% a TeX un margen de estirado de emergencia antes de desbordar la línea
% (evita Overfull \hbox menores sin tocar el contenido).
\setlength{\emergencystretch}{1em}

% --- Datos globales y macros de la memoria (única fuente: no duplicar) ---
% =========================================================================
% Variables globales del documento.
% Cambiar un dato AQUÍ lo actualiza en todo el documento (portada,
% metadatos del PDF...). No escribir nunca estos datos "a mano" en otros
% archivos: usar siempre estos comandos.
% =========================================================================

% --- Trabajo ---
\newcommand{\thesisTitle}{BUKIT: red social orientada a objetivos}% TODO: título definitivo
\newcommand{\thesisSubtitle}{Memoria del Trabajo de Fin de Grado}%
\newcommand{\thesisDegree}{Grado en Ingeniería Informática de Gestión y Sistemas de Información}%
\newcommand{\thesisCourse}{2026--2027}% curso académico
\newcommand{\thesisDate}{Bilbao, \today}% TODO: <Lugar, día, mes, año> de la entrega/defensa
\newcommand{\thesisYear}{2027}% año de elaboración y defensa

% --- Personas (la plantilla del centro usa "Apellido1 Apellido2, Nombre") ---
\newcommand{\thesisAuthor}{Liviu Deleanu}% nombre natural (para metadatos del PDF)
\newcommand{\thesisAuthorFormal}{Deleanu, Liviu}% TODO: confirmar formato con dos apellidos
\newcommand{\thesisDirector}{Villamañe Gironés, Mikel}%
\newcommand{\thesisCoDirector}{Apellido1 Apellido2, Nombre}% TODO: solo si hay codirector/a

% --- Universidad ---
\newcommand{\thesisUniversity}{Universidad del País Vasco / Euskal Herriko Unibertsitatea}%
\newcommand{\thesisFaculty}{Escuela de Ingeniería de Bilbao}%

% --- Palabras clave del resumen trilingüe ---
\newcommand{\thesisKeywordsEs}{red social, objetivos, monorepo, TypeScript, PostgreSQL}% TODO
\newcommand{\thesisKeywordsEu}{sare soziala, helburuak, monorepo, TypeScript, PostgreSQL}% TODO: revisar euskera
\newcommand{\thesisKeywordsEn}{social network, goals, monorepo, TypeScript, PostgreSQL}% TODO

% =========================================================================
% Comandos propios. Definir aquí cualquier cosa que se repita en el texto:
% si mañana cambia cómo se escribe, se cambia en un único sitio.
% =========================================================================

% Nombre del producto, siempre con el mismo formato.
\newcommand{\bukit}{\textsc{Bukit}\xspace}

% Términos en inglés, en cursiva de forma consistente: \eng{backend}
\newcommand{\eng}[1]{\textit{#1}}

% Notas de trabajo, visibles en rojo en el PDF. Buscar "todonote" y
% eliminarlas todas antes de entregar.
\newcommand{\todonote}[1]{\textcolor{red}{\textbf{[TODO: #1]}}}


% --- Cabecera y pie (diseño de referencia, sin logos) ---
% La caja de cabecera replica la del .tex de referencia (celda ancha de
% contexto + celda estrecha de identidad). La celda derecha llevaba el logo
% del centro: aquí va texto porque los logos UPV/EHU siguen pendientes de
% la fase de diseño (ver style.tex); cuando existan, sustituir por
% \includegraphics. 11 + 4,9 = 15,9 cm < 16 cm de texto: sin desbordes.
\usepackage{fancyhdr}
\newcommand{\riesgosHeader}{%
  \setlength{\tabcolsep}{0pt}%
  \renewcommand{\arraystretch}{1}%
  \begin{tabular}{|p{11cm}|p{4.9cm}|}
    \hline
    \parbox[c][2cm][c]{11cm}{\centering\bfseries
      {\fontsize{16}{19}\selectfont TRABAJO DE FIN DE GRADO\\[4pt]Curso \thesisCourse}%
    } &
    \parbox[c][2cm][c]{4.9cm}{\centering
      {\fontsize{14}{17}\selectfont\bfseries\bukit}\\[3pt]
      {\small\thesisFaculty}%
    }\\
    \hline
  \end{tabular}%
}
\pagestyle{fancy}
\fancyhf{}
\renewcommand{\headrulewidth}{0pt}
\fancyhead[C]{\riesgosHeader}
\fancyfoot[C]{\thepage}
% La primera página usa estilo "plain": se redefine para que también lleve
% la cabecera (mismo truco que el .tex de referencia).
\fancypagestyle{plain}{%
  \fancyhf{}%
  \renewcommand{\headrulewidth}{0pt}%
  \fancyhead[C]{\riesgosHeader}%
  \fancyfoot[C]{\thepage}%
}

% --- Símbolos y campos de formulario del diseño de referencia ---
\newcommand{\fieldline}[1]{\underline{\hspace{#1}}}
\newcommand{\bigline}{\noindent\rule{\textwidth}{0.4pt}}
\newcommand{\checkbox}{\(\square\)}
\newcommand{\checkedbox}{\(\blacksquare\)}
% Símbolo ⟳ = "si se rellena este apartado, nueva versión del plan".
% El original lo resolvía con un SVG externo (paquete svg + Inkscape);
% aquí basta el símbolo de amssymb: cero dependencias nuevas.
\newcommand{\updatesymbol}{\raisebox{-0.2ex}{\(\circlearrowright\)}}
% Resaltado amarillo del identificador de riesgo en los títulos.
\newcommand{\markfield}[1]{\colorbox{yellow!35}{\strut #1}}

% --- Enlaces y metadatos del PDF (al final, como exige hyperref) ---
% El pdftitle lo fija cada documento con \hypersetup tras cargar esto.
\usepackage[hidelinks]{hyperref}
\hypersetup{pdfauthor={\thesisAuthor}}

\hypersetup{pdftitle={Seguimiento de riesgos — 25-10-2026 (S9)}}

\begin{document}

\begin{center}
  {\LARGE\bfseries Seguimiento de riesgos — S9}\\[6pt]
\end{center}

\textbf{Fecha:} 25-10-2026 · \textbf{Sprint:} S9 (19-10 \(\to\) 25-10-2026, tercer sprint de F2) · \textbf{Responsable:} \thesisAuthor

Regla de revisión: \textbf{críticos} cada sprint; \textbf{significativos} cada 2 sprints; \textbf{moderados} cada 4 sprints; cualquier riesgo, \textbf{de inmediato}, si se observa su indicador clave.
Una semana declarada a cero (cap.~10 de la memoria) no cierra sprint y no tiene seguimiento: lo que tocara por cadencia se revisa en el siguiente.

\textbf{El símbolo \updatesymbol\ indica que, si se rellena ese apartado, hay que emitir una nueva versión completa del plan de riesgos} en \texttt{../plan-riesgos/} y registrar el cambio en la tabla de seguimiento del cap.~7 de la memoria.

Riesgos revisados:
\textbf{R1--R7 y R10 por cadencia} (críticos, semanal);
\textbf{R8, R9, R11, R14 y R15 por cadencia} (significativos: les tocaba en S9);
y \textbf{R12 por disparador}: su indicador saltó el 19-10.
R13 (moderado) toca en S11.
Próximas revisiones por cadencia: significativos en S11 y moderados (R12 y R13) en S11; R12 se revisa además cada sprint mientras siga activo.

\textbf{Plan emitido a partir de este seguimiento:} ninguno.
El único cambio es el estado de R12, que pasa a activo sin cambio de probabilidad, impacto ni medidas; desde este seguimiento, un cambio así se registra en la tabla del cap.~7 y aquí, sin reemitir el plan (decisión del autor, ver «Reflexión»).

Datos del sprint: velocidad de S9 = \textbf{3 SP} (US-08, cerrada el 21-10 y subida el 22-10, con la que arranca EP03); horas reales según el diario, 3~h~20 (20~min el 21-10 y 3~h el 25-10), frente a las 5~h declaradas: 67~\%.
Dos sesiones (miércoles y domingo).
La historia la implementó el asistente en una sesión autónoma de 20 minutos del autor, con las decisiones (enrutador, estructura por funcionalidad, puerto fijo) tomadas por él y cuatro rondas de revisión; las tres horas del domingo fueron del cierre del sprint y la lectura de la memoria.
Equivalencia observada: no separable ni representativa (ver la retrospectiva del diario); se mantiene la de planificación, 2,0~h/SP, y desde S10 el diario desglosa las horas de la historia y las del resto, para recalibrar en S12 con datos limpios.

\bigline

\section{Medidas de mitigación (prevención)}

Para cada medida de los riesgos que tocan revisión en este seguimiento: ¿se está aplicando?
Si no, justificar y decidir si se mantiene/modifica/elimina.

% Bloques de 3-4 filas (tabularx no parte entre páginas; receta del S2).
\Needspace*{24\baselineskip}
\underline{\textbf{Riesgos críticos (R1--R4)}}

\begin{tabularx}{\textwidth}{|C{0.9cm}|Y|Y|C{2.2cm}|L{2.4cm}|}
\hline
\multicolumn{1}{|C{0.9cm}|}{\textbf{Id}} &
\multicolumn{1}{>{\centering\arraybackslash}X|}{\textbf{Medida de mitigación}} &
\multicolumn{1}{>{\centering\arraybackslash}X|}{\textbf{¿Aplicada?}} &
\multicolumn{1}{C{2.2cm}|}{\textbf{Justificación (del NO)}} &
\multicolumn{1}{C{2.4cm}|}{\textbf{Decisión}} \\
\hline
R1 & Toda idea nueva entra por el backlog. & \checkedbox\ Sí (EP18 y la US-19 entraron por el backlog en la planificación del 21-10, con su requisito RNF-09 y su anclaje en el cap.~3; el sprint comprometió solo la US-08).\newline \checkbox\ No & — & \checkedbox\ Mantener\newline \checkbox\ Modificar\newline \checkbox\ Eliminar \\
\hline
R2 & Medir velocidad y recalibrar en cada cierre. & \checkedbox\ Sí (velocidad medida, 3 SP; la recalibración se debatió en la retro y se aplaza a S12 con el desglose de horas, ver «Datos del sprint»).\newline \checkbox\ No & — & \checkedbox\ Mantener\newline \checkbox\ Modificar\newline \checkbox\ Eliminar \\
\hline
R3 & Diario al cierre de cada sesión; memoria con hueco; ninguna historia cerrada sin documentar; auditoría semanal de coherencia. & \checkedbox\ Sí (caps.~6, 9 y 10 actualizados en la sesión de la US-08; el backlog nuevo en los caps.~3, 5, 10, 13 y 14 ese mismo día; informe de coherencia del S8 aplicado por el autor).\newline \checkbox\ No & — & \checkedbox\ Mantener\newline \checkbox\ Modificar\newline \checkbox\ Eliminar \\
\hline
R4 & Citar al escribir; pendientes con nota visible. & \checkedbox\ Sí (cuatro fuentes nuevas —React Router, TanStack Router, GSAP y Lenis— citadas al escribir el cap.~6).\newline \checkbox\ No & — & \checkedbox\ Mantener\newline \checkbox\ Modificar\newline \checkbox\ Eliminar \\
\hline
\end{tabularx}

\Needspace*{26\baselineskip}
\underline{\textbf{Riesgos críticos (R5--R7 y R10)}}

\begin{tabularx}{\textwidth}{|C{0.9cm}|Y|Y|C{2.2cm}|L{2.4cm}|}
\hline
\multicolumn{1}{|C{0.9cm}|}{\textbf{Id}} &
\multicolumn{1}{>{\centering\arraybackslash}X|}{\textbf{Medida de mitigación}} &
\multicolumn{1}{>{\centering\arraybackslash}X|}{\textbf{¿Aplicada?}} &
\multicolumn{1}{C{2.2cm}|}{\textbf{Justificación (del NO)}} &
\multicolumn{1}{C{2.4cm}|}{\textbf{Decisión}} \\
\hline
R5 & Revisar todo antes de integrar; nada sin comprobaciones. & \checkbox\ Sí\newline \checkedbox\ No, en parte: el código de la US-08 pasó cuatro rondas de revisión (estructura, código, pruebas y comentarios) por agentes y por el asistente, con lint, tipos y pruebas en verde, pero el autor aún no lo ha leído. & La sesión fue autónoma por decisión del autor, que fijó las decisiones (enrutador, estructura, puerto) y revisó el plan del backlog, no el código. & \checkedbox\ Mantener, con acción: el autor lee \texttt{apps/web/src} (doce archivos) antes de la US-19\newline \checkbox\ Modificar\newline \checkbox\ Eliminar \\
\hline
R6 & Investigación acotada ante tecnología nueva. & \checkedbox\ Sí (React Router 8 y Vite 8, nuevos: comparativa con TanStack Router el 21-10 y la historia cerrada en su sprint).\newline \checkbox\ No & — & \checkedbox\ Mantener\newline \checkbox\ Modificar\newline \checkbox\ Eliminar \\
\hline
R7 & Minimización desde el diseño (RGPD). & \checkedbox\ Sí (la web no maneja datos personales todavía y el modelo de datos no cambió).\newline \checkbox\ No & — & \checkedbox\ Mantener\newline \checkbox\ Modificar\newline \checkbox\ Eliminar \\
\hline
R10 & Margen deliberado + documentación continua. & \checkedbox\ Sí (sin parón: sesiones el miércoles y el domingo).\newline \checkbox\ No & — & \checkedbox\ Mantener\newline \checkbox\ Modificar\newline \checkbox\ Eliminar \\
\hline
\end{tabularx}

\Needspace*{22\baselineskip}
\underline{\textbf{Riesgos significativos (R8, R9 y R11)}}

\begin{tabularx}{\textwidth}{|C{0.9cm}|Y|Y|C{2.2cm}|L{2.4cm}|}
\hline
\multicolumn{1}{|C{0.9cm}|}{\textbf{Id}} &
\multicolumn{1}{>{\centering\arraybackslash}X|}{\textbf{Medida de mitigación}} &
\multicolumn{1}{>{\centering\arraybackslash}X|}{\textbf{¿Aplicada?}} &
\multicolumn{1}{C{2.2cm}|}{\textbf{Justificación (del NO)}} &
\multicolumn{1}{C{2.4cm}|}{\textbf{Decisión}} \\
\hline
R8 & Horas por sesión en el diario; cada retro compara; semanas a cero declaradas antes. & \checkedbox\ Sí (horas registradas en cada sesión; ninguna semana a cero que declarar).\newline \checkbox\ No & — & \checkedbox\ Mantener\newline \checkbox\ Modificar\newline \checkbox\ Eliminar \\
\hline
R9 & Objetivos secundarios o versiones simuladas; capas gratuitas. & \checkedbox\ Sí (Render y Neon en capa gratuita, sin incidencias ni anuncios de cambios; la web no depende de ningún servicio externo).\newline \checkbox\ No & — & \checkedbox\ Mantener\newline \checkbox\ Modificar\newline \checkbox\ Eliminar \\
\hline
R11 & Todo en git con remoto y subidas frecuentes. & \checkedbox\ Sí (cada sesión de la semana cerró con su subida: 22-10 y 25-10; sin anomalías del equipo).\newline \checkbox\ No & — & \checkedbox\ Mantener\newline \checkbox\ Modificar\newline \checkbox\ Eliminar \\
\hline
\end{tabularx}

\Needspace*{30\baselineskip}
\underline{\textbf{Riesgos significativos R14 y R15, y moderado R12 (por disparador)}}

\begin{tabularx}{\textwidth}{|C{0.9cm}|Y|Y|C{2.2cm}|L{2.4cm}|}
\hline
\multicolumn{1}{|C{0.9cm}|}{\textbf{Id}} &
\multicolumn{1}{>{\centering\arraybackslash}X|}{\textbf{Medida de mitigación}} &
\multicolumn{1}{>{\centering\arraybackslash}X|}{\textbf{¿Aplicada?}} &
\multicolumn{1}{C{2.2cm}|}{\textbf{Justificación (del NO)}} &
\multicolumn{1}{C{2.4cm}|}{\textbf{Decisión}} \\
\hline
R14 & Auditoría en la CI + Dependabot con triaje semanal y alertas + casos de abuso en la Definición de Hecho. & \checkedbox\ Sí, y esta semana demostró su valor: el \texttt{pnpm audit} del triaje encontró dos altas de producción que ni Dependabot proponía corregir ni GitHub mostraba como alertas; resueltas el mismo día (ver ficha). La US-08 no toca autenticación ni datos personales.\newline \checkbox\ No & — & \checkedbox\ Mantener\newline \checkbox\ Modificar\newline \checkbox\ Eliminar \\
\hline
R15 & Dimensionar los agentes antes de automatizar; modelos de coste contenido fijados; tareas pesadas de noche. & \checkedbox\ Sí (las cuatro rondas de revisión de la US-08 fueron con un agente cada vez, en secuencia, sin agotar la cuota; lección del 15-10 aplicada).\newline \checkbox\ No & — & \checkedbox\ Mantener\newline \checkbox\ Modificar\newline \checkbox\ Eliminar \\
\hline
R12 & Capítulos a revisión conforme se cierran; versión vigente cuando haya cambios que merezcan revisión, según la disponibilidad de ambos. & \checkedbox\ Sí (la medida nueva se aplica: sin entregas forzadas; el trabajo sigue sin depender de la revisión).\newline \checkbox\ No & — & \checkedbox\ Mantener (el indicador saltó: ver ficha)\newline \checkbox\ Modificar\newline \checkbox\ Eliminar \\
\hline
\end{tabularx}

\Needspace*{8\baselineskip}
\underline{\textbf{Cambios realizados en las medidas de mitigación}} \updatesymbol

\begin{tabularx}{\textwidth}{|C{1.4cm}|Y|Y|Y|}
\hline
\multicolumn{1}{|C{1.4cm}|}{\textbf{Id}} &
\multicolumn{1}{>{\centering\arraybackslash}X|}{\textbf{Medida original}} &
\multicolumn{1}{>{\centering\arraybackslash}X|}{\textbf{Cambio realizado}} &
\multicolumn{1}{>{\centering\arraybackslash}X|}{\textbf{Justificación}} \\
\hline
— & Ninguno en este seguimiento. & — & — \\
\hline
\end{tabularx}

\Needspace*{8\baselineskip}
\underline{\textbf{Nuevas medidas de mitigación a añadir}} \updatesymbol

\begin{tabularx}{\textwidth}{|C{1.4cm}|Y|Y|}
\hline
\multicolumn{1}{|C{1.4cm}|}{\textbf{Id}} &
\multicolumn{1}{>{\centering\arraybackslash}X|}{\textbf{Nueva medida}} &
\multicolumn{1}{>{\centering\arraybackslash}X|}{\textbf{Justificación}} \\
\hline
— & Ninguna en este seguimiento. & — \\
\hline
\end{tabularx}

\bigline

\section{Riesgos a revisar en este seguimiento}

Los riesgos \textbf{sin novedad y sin observación relevante} se agrupan en la tabla siguiente (indicador en su valor esperado, sin salto, estado «sin novedad»); los que dejaron señal en el periodo (R12 y R14) conservan su ficha completa a continuación.

\Needspace*{30\baselineskip}
\begin{tabularx}{\textwidth}{|C{0.9cm}|Y|C{2.6cm}|}
\hline
\multicolumn{1}{|C{0.9cm}|}{\textbf{Id}} &
\multicolumn{1}{>{\centering\arraybackslash}X|}{\textbf{Indicador: valor real (esperado 0 en todos)}} &
\multicolumn{1}{C{2.6cm}|}{\textbf{Nuevo estado}} \\
\hline
R1 & 0 — el sprint no ganó puntos tras la planificación; EP18 (deseable, 8 SP) y la US-19 (3 SP, comprometida) entraron por el backlog con su requisito y su anclaje en el cap.~3, como pide la medida. & Sin novedad \\
\hline
R2 & Indicador (tres primeros sprints) ya evaluado en S3. S9: 3 SP; la equivalencia no se recalibra por falta de horas separables, y el diario las desglosa desde S10. & Sin novedad \\
\hline
R3 & 0 — capítulos al día en la sesión del cambio; diario escrito en sus dos sesiones. & Sin novedad \\
\hline
R4 & 0 — cuatro fuentes nuevas, citadas al escribir. & Sin novedad \\
\hline
R5 & Observación: las decisiones de la US-08 están en el diario con sus alternativas y las tomó el autor; el código generado se integró revisado por agentes y pruebas, no por el autor. Acción: leer \texttt{apps/web/src} antes de la US-19; si no pudiera explicarlo, el indicador habría saltado. & Sin novedad \\
\hline
R6 & 0 — la US-08 (3 SP, 6~h previstas) se cerró en su sprint sin desvío. & Sin novedad \\
\hline
R7 & 0 — el modelo de datos no cambió; la web no maneja datos personales. & Sin novedad \\
\hline
R8 & Observación: S9 por debajo del 70~\% de las 5~h declaradas, tras un S8 al 170~\%; una semana, no dos consecutivas, y con causa identificada (sesión autónoma del asistente para la US-08). & Sin novedad \\
\hline
R9 & 0 — sin anuncios de precios ni caídas en Render ni en Neon. & Sin novedad \\
\hline
R10 & 0 — sin parón; dos sesiones en la semana. & Sin novedad \\
\hline
R11 & 0 — cada sesión subió sus cambios; sin anomalías del equipo. & Sin novedad \\
\hline
R15 & 0 — la cuota semanal no se agotó; un flujo de agentes cada vez. & Sin novedad \\
\hline
\end{tabularx}

% ------------------ R12 ------------------
\Needspace*{30\baselineskip}
\subsection*{Riesgo: \markfield{R12 — Ciclos de revisión con el tutor} (por disparador)}

\textbf{Indicadores (clave y de resolución)}

\begin{tabularx}{\textwidth}{|Z|Z|Z|Z|Z|}
\hline
\textbf{Indicador} & \textbf{Valor esperado} & \textbf{Valor real} & \textbf{¿Cumplido?} & \textbf{Acción / observaciones} \\
\hline
Consulta/capítulo >2 semanas sin respuesta, o revisión con cambios estructurales sobre capítulo cerrado. & 0 & \textbf{Sí, el 19-10}: la entrega del 05-10 (memoria completa y cuatro consultas) cumplió dos semanas sin respuesta del tutor; a 25-10, veinte días. & \checkedbox\ Sí\quad \checkbox\ No & Previsto en el S8, que mantuvo el indicador a sabiendas. La contingencia se aplica: el trabajo no depende de la revisión (US-08 y backlog siguieron). Sin recordatorio, por el acuerdo del 15-10 de entregas según disponibilidad. Se revisa cada sprint mientras siga activo; la consulta del cap.~3 (§3.2 frente a §3.5.1) sigue abierta. \\
\hline
\end{tabularx}

\textbf{Nueva probabilidad (1–3):} \underline{\ 1\ } · \textbf{Nuevo impacto (1–3):} \underline{\ 2\ } · \textbf{Nueva severidad:} \checkbox\ Crítico · \checkbox\ Significativo · \checkedbox\ Moderado\ \updatesymbol\ (sin cambio: una sola observación; se reevalúa en S11 si persiste)

\textbf{¿Cerrar el riesgo?} \checkbox\ Sí · \checkedbox\ No\ \updatesymbol

\textbf{¿Se ha materializado?} \checkbox\ Sí · \checkedbox\ No\ \updatesymbol

\textbf{Nuevo estado} (cap.~7): \checkbox\ Identificado · \checkbox\ Sin novedad · \checkedbox\ Activo · \checkbox\ Materializado · \checkbox\ Resuelto · \checkbox\ Cerrado

\Needspace*{10\baselineskip}
\textbf{Plan de contingencia (solo si se ha materializado)}: no aplicable; la demora todavía no retrasa ningún cierre.
El cambio de estado se registra en la tabla del cap.~7 sin reemitir el plan (ver «Reflexión»).

% ------------------ R14 ------------------
\Needspace*{30\baselineskip}
\subsection*{Riesgo: \markfield{R14 — Verificación de seguridad insuficiente}}

\textbf{Indicadores (clave y de resolución)}

\begin{tabularx}{\textwidth}{|Z|Z|Z|Z|Z|}
\hline
\textbf{Indicador} & \textbf{Valor esperado} & \textbf{Valor real} & \textbf{¿Cumplido?} & \textbf{Acción / observaciones} \\
\hline
Historia con autenticación, autorización o datos personales cerrada sin casos de abuso, o vulnerabilidad alta sin resolver: las de producción, dentro del sprint; las de desarrollo, al cierre del siguiente. & 0 & 0 — dos altas de producción aparecieron durante la semana (zod anterior a la 4.6.6, agotamiento de memoria con colecciones acotadas; body-parser anterior a la 2.4.0, denegación de servicio con cuerpos comprimidos) y se resolvieron el 25-10, dentro del sprint: migración a Zod 4.6.7 y lockfile refrescado. \texttt{pnpm audit} a cero en producción y en desarrollo. & \checkbox\ Sí\quad \checkedbox\ No & Las cazó el \texttt{pnpm audit} semanal del triaje: Dependabot proponía subir zod a la 3.25.76, todavía vulnerable, y GitHub no mostraba alertas abiertas. Zod 4 se había aplazado a F4 el 11-10 por no aportar a las historias; la seguridad lo adelantó, con una migración pequeña (correo normalizado y después \texttt{z.email()}, \texttt{z.uuid()}), 81 pruebas en verde. \\
\hline
\end{tabularx}

\textbf{Nueva probabilidad (1–3):} \fieldline{0.8cm} · \textbf{Nuevo impacto (1–3):} \fieldline{0.8cm} · \textbf{Nueva severidad:} \checkbox\ Crítico · \checkbox\ Significativo · \checkbox\ Moderado\ \updatesymbol\ (sin reevaluación: el indicador no ha saltado)

\textbf{¿Cerrar el riesgo?} \checkbox\ Sí · \checkedbox\ No\ \updatesymbol

\textbf{¿Se ha materializado?} \checkbox\ Sí · \checkedbox\ No\ \updatesymbol

\textbf{Nuevo estado} (cap.~7): \checkbox\ Identificado · \checkedbox\ Sin novedad · \checkbox\ Activo · \checkbox\ Materializado · \checkbox\ Resuelto · \checkbox\ Cerrado

\textbf{Plan de contingencia (solo si se ha materializado)}: no aplicable.

\Needspace*{10\baselineskip}
\bigline

\section{Reflexión desde el último seguimiento}

\textbf{¿Ha ocurrido algo reseñable desde el último seguimiento?} \checkedbox\ Sí · \checkbox\ No

Una semana desde el seguimiento del S8.
La US-08 abrió EP03 con el cliente web organizado por funcionalidad, y el backlog incorporó EP18, la US-19 y la preparación de la defensa (caps.~3, 5, 10, 13 y 14).
Dos vulnerabilidades altas de producción obligaron a adelantar Zod 4, aplazado en S7.
R12 saltó el 19-10, como preveía el S8.
\textbf{Decisión del autor sobre el versionado del plan}: un cambio de estado que no altera probabilidad, impacto, medidas ni contingencias (la activación de R12) se registra en la tabla del cap.~7 y en el seguimiento, sin reemitir el plan; la columna de estado de cada versión refleja la situación al emitirla.
Alternativa descartada: una v8 solo para cambiar «Identificado» por «Activo» en R12; porqué: siete versiones en catorce semanas, dos de ellas por redacción, y el objetivo fijado el 17-10 de reemitir el plan solo cuando cambie lo que el plan decide.
La regla se añade al cap.~7 y a los README del plan y del seguimiento.

\Needspace*{6\baselineskip}
\textbf{Nuevos riesgos identificados} \updatesymbol

\begin{tabularx}{\textwidth}{|C{1.4cm}|Y|C{0.8cm}|C{0.8cm}|C{1cm}|C{2.4cm}|}
\hline
\multicolumn{1}{|C{1.4cm}|}{\textbf{Id}} &
\multicolumn{1}{>{\centering\arraybackslash}X|}{\textbf{Descripción}} &
\multicolumn{1}{C{0.8cm}|}{\textbf{P}} &
\multicolumn{1}{C{0.8cm}|}{\textbf{I}} &
\multicolumn{1}{C{1cm}|}{\textbf{Exp.}} &
\multicolumn{1}{C{2.4cm}|}{\textbf{Severidad}} \\
\hline
— & Ninguno. & — & — & — & — \\
\hline
\end{tabularx}

\textbf{¿Afecta algo a los objetivos o al alcance (cap.~3 de la memoria)?} \checkbox\ Sí · \checkedbox\ No

\textbf{Justificación:} EP18 entra como deseable, sin criterio de éxito, y la US-19 como comprometida dentro de EP03 (11 SP); F2 proyecta 44 SP frente a 40 de capacidad, diferencia que asume la holgura del plan si la recalibración de S12 lo confirma, y el alcance del cap.~3 no cambia.

\textbf{Otros eventos relevantes:} el autor leyó la memoria completa (pendiente desde el 04-10) y declaró a cero la semana siguiente (S10), por exámenes, antes de empezarla, como pide la regla del cap.~10: no contará como parón (R10) ni como falta de horas (R8); la CI pasó a Ubuntu 26 sin cambios (primera ejecución el 22-10, R13); TypeScript 7 evaluada y aplazada (typescript-eslint admite hasta la 6.0 y la 7 retira \texttt{baseUrl}); triaje de Dependabot del 25-10 con sus cinco propuestas (el tope): cuatro fusionadas de una en una con su CI en verde sobre el \texttt{main} nuevo (vitest 5.0.5, globals 17.13.0 y las mayores eslint-config-prettier 10.1.8 y ESLint 10.13.0, admitidas por los pares de sus complementos), y la de zod 3.25.76 (\#29) cerrada al quedar superada por la migración.

\end{document}
